% GENERATED FILE. Edit canonical lessons, then run npm run generate:books.
% canonical-source-hash: fnv1a64:0fb677d4ef76198e
% canonical-lessons: ZH-C275-daoqian, ZH-C275-yuanliang, ZH-C275-daying, ZH-C275-jieshou, ZH-C275-suibian

\chapter{Saying Sorry}
\label{ch:zh-saying-sorry}
\hlchaptermodality{275}

\begin{chapteropening}
\textbf{By the end of this chapter:} \emph{I can apologise, forgive, say yes to a request, accept an apology, and say it is up to you.}

Complete the last lesson of chapter 275: I can apologise, forgive, say yes to a request, accept an apology, and say it is up to you.
\end{chapteropening}

\section[dàoqiàn]{dàoqiàn — to apologise}

\label{lesson:ZH-C275-daoqian}



\begin{quote}
Before the new one: say the Mandarin for handicrafts, then the Mandarin for a sports fan.
\end{quote}

\subsection*{You'll want to know: dàoqiàn}
\emph{dàoqiàn} — \textquotedblleft{}to apologise\textquotedblright{}.

\begin{quote}
\emph{tā dàoqiàn le} — He apologised.
\end{quote}

Its characters come later; for now it is a word for the ear.

\begin{cousinweb}
\emph{dàoqiàn} is built from \emph{dào} and \emph{qiàn}, both inherited from Middle Chinese, \emph{dawX} and \emph{kheamH} (Guangyun, in Baxter's transcription), which glosses the first 'a road' and 'a principle'. It gives the second other senses; how the pair came to mean an apology is not settled.
\end{cousinweb}

\begin{grammarlens}[title={What you've built}]
One more word for saying no and asking.
\end{grammarlens}

\subsection*{Your turn}
\begin{itemize}
\raggedright
  \item \emph{Say it:} \emph{dàoqiàn}
  \item \emph{Say it:} \emph{dàoqiàn}, once more
  \item \emph{Say it:} \emph{tā dàoqiàn le} — He apologised.
  \item \emph{Recall:} say the Mandarin for handicrafts, then the Mandarin for a sports fan, then say \emph{dàoqiàn} again
\end{itemize}

\begin{tcolorbox}[breakable,colback=teal!4,colframe=teal!35!black,title={Before you move on}]
What does dàoqiàn mean? (\textquotedblleft{}To apologise\textquotedblright{}.) Say it once more.
\end{tcolorbox}
\section[yuánliàng]{yuánliàng — to forgive}

\label{lesson:ZH-C275-yuanliang}



\begin{quote}
Before the new one: say the Mandarin for a sports fan, then the Mandarin for to apologise.
\end{quote}

\subsection*{You'll want to know: yuánliàng}
\emph{yuánliàng} — \textquotedblleft{}to forgive\textquotedblright{}.

\begin{quote}
\emph{qǐng yuánliàng wǒ} — Please forgive me.
\end{quote}

Its characters come later; for now it is a word for the ear.

\begin{cousinweb}
\emph{yuánliàng} is built from \emph{yuán} and \emph{liàng} 'to trust', both inherited from Middle Chinese, \emph{ngjwon} and \emph{ljangH} (Guangyun, in Baxter's transcription), which glosses the first 'a broad plain' and the second 'to trust'. How the pair came to mean forgiving is not settled.
\end{cousinweb}

\begin{grammarlens}[title={What you've built}]
One more word for saying no and asking.
\end{grammarlens}

\subsection*{Your turn}
\begin{itemize}
\raggedright
  \item \emph{Say it:} \emph{yuánliàng}
  \item \emph{Say it:} \emph{yuánliàng}, once more
  \item \emph{Say it:} \emph{qǐng yuánliàng wǒ} — Please forgive me.
  \item \emph{Recall:} say the Mandarin for a sports fan, then the Mandarin for to apologise, then say \emph{yuánliàng} again
\end{itemize}

\begin{tcolorbox}[breakable,colback=teal!4,colframe=teal!35!black,title={Before you move on}]
What does yuánliàng mean? (\textquotedblleft{}To forgive\textquotedblright{}.) Say it once more.
\end{tcolorbox}
\section[dāying]{dāying — to agree (to a request); to say yes}

\label{lesson:ZH-C275-daying}



\begin{quote}
Before the new one: say the Mandarin for to apologise, then the Mandarin for to forgive.
\end{quote}

\subsection*{You'll want to know: dāying}
\emph{dāying} — \textquotedblleft{}to agree (to a request); to say yes\textquotedblright{}.

\begin{quote}
\emph{tā dāying le} — He said yes.
\end{quote}

Its characters come later; for now it is a word for the ear.

\begin{cousinweb}
\emph{dāying} is built from \emph{dā} 'to answer' and \emph{yìng} 'to respond', both inherited from Middle Chinese, \emph{top} and \emph{'ingH} (Guangyun, in Baxter's transcription), which glosses the first 'to match' and the second 'things answering each other'. In the word the second loses its tone.
\end{cousinweb}

\begin{grammarlens}[title={What you've built}]
One more word for saying no and asking.
\end{grammarlens}

\subsection*{Your turn}
\begin{itemize}
\raggedright
  \item \emph{Say it:} \emph{dāying}
  \item \emph{Say it:} \emph{dāying}, once more
  \item \emph{Say it:} \emph{tā dāying le} — He said yes.
  \item \emph{Recall:} say the Mandarin for to apologise, then the Mandarin for to forgive, then say \emph{dāying} again
\end{itemize}

\begin{tcolorbox}[breakable,colback=teal!4,colframe=teal!35!black,title={Before you move on}]
What does dāying mean? (\textquotedblleft{}To agree (to a request); to say yes\textquotedblright{}.) Say it once more.
\end{tcolorbox}
\section[jiēshòu]{jiēshòu — to accept}

\label{lesson:ZH-C275-jieshou}



\begin{quote}
Before the new one: say the Mandarin for to forgive, then the Mandarin for to say yes to a request.
\end{quote}

\subsection*{You'll want to know: jiēshòu}
\emph{jiēshòu} — \textquotedblleft{}to accept\textquotedblright{}.

\begin{quote}
\emph{wǒ jiēshòu nǐ de dàoqiàn} — I accept your apology.
\end{quote}

Its characters come later; for now it is a word for the ear.

\begin{cousinweb}
\emph{jiēshòu} is built from \emph{jiē} and \emph{shòu}, both inherited from Middle Chinese, \emph{tsjep} and \emph{dzyuwX} (Guangyun, in Baxter's transcription), which glosses the first 'to join' and 'to hold', and the second 'to take in'.
\end{cousinweb}

\begin{grammarlens}[title={What you've built}]
One more word for saying no and asking.
\end{grammarlens}

\subsection*{Your turn}
\begin{itemize}
\raggedright
  \item \emph{Say it:} \emph{jiēshòu}
  \item \emph{Say it:} \emph{jiēshòu}, once more
  \item \emph{Say it:} \emph{wǒ jiēshòu nǐ de dàoqiàn} — I accept your apology.
  \item \emph{Recall:} say the Mandarin for to forgive, then the Mandarin for to say yes to a request, then say \emph{jiēshòu} again
\end{itemize}

\begin{tcolorbox}[breakable,colback=teal!4,colframe=teal!35!black,title={Before you move on}]
What does jiēshòu mean? (\textquotedblleft{}To accept\textquotedblright{}.) Say it once more.
\end{tcolorbox}
\section[suíbiàn]{suíbiàn — as you like; casual}

\label{lesson:ZH-C275-suibian}



\begin{quote}
Before the new one: say the Mandarin for to say yes to a request, then the Mandarin for to accept.
\end{quote}

\subsection*{You'll want to know: suíbiàn}
\emph{suíbiàn} — \textquotedblleft{}as you like; casual\textquotedblright{}.

\begin{quote}
\emph{nǐ suíbiàn chī} — Help yourself to anything.
\end{quote}

Its characters come later; for now it is a word for the ear.

\begin{cousinweb}
\emph{suíbiàn} is built from \emph{suí} 'to follow' and \emph{biàn} 'convenient', both inherited from Middle Chinese, \emph{zjwe} and \emph{bjienH} (Guangyun, in Baxter's transcription), which glosses the first 'to follow' and the second 'to benefit'.
\end{cousinweb}

\begin{grammarlens}[title={What you've built}]
One more word for saying no and asking.
\end{grammarlens}

\subsection*{Your turn}
\begin{itemize}
\raggedright
  \item \emph{Say it:} \emph{suíbiàn}
  \item \emph{Say it:} \emph{suíbiàn}, once more
  \item \emph{Say it:} \emph{nǐ suíbiàn chī} — Help yourself to anything.
  \item \emph{Recall:} say the Mandarin for to say yes to a request, then the Mandarin for to accept, then say \emph{suíbiàn} again
\end{itemize}

\begin{tcolorbox}[breakable,colback=teal!4,colframe=teal!35!black,title={Before you move on}]
What does suíbiàn mean? (\textquotedblleft{}As you like; casual\textquotedblright{}.) Say it once more.
\end{tcolorbox}
